\section{Discussion}

\subsection{Effet de l'heuristique}

Sur le 8-Puzzle moyen, les trois variantes d'\AStar{} trouvent une solution de
coût 28. La qualité de l'heuristique modifie toutefois fortement l'effort :
Manhattan développe 1\,440 noeuds, contre 59\,067 avec les tuiles mal placées
et 177\,134 avec $h=0$. Manhattan réduit ainsi le nombre de noeuds développés
de 99,2\,\% par rapport à la recherche sans information heuristique. Dans
cette expérience, le temps mesuré diminue également de 1,388574\,s à
0,014178\,s. Les temps étant mesurés une seule fois, cette dernière
comparaison reste indicative.

Greedy développe seulement 385 noeuds avec Manhattan, mais produit une
solution de 48 coups, alors que \AStar{} en trouve une de 28 coups. La
réduction du travail de recherche ne garantit donc pas l'optimalité de la
solution.

\subsection{Tree-Search contre Graph-Search}

Sur le labyrinthe pondéré, le nombre de noeuds expansés par \UCS{} passe de
596 en \TreeSearch{} à 15 en \GraphSearch{}; pour \BFS{}, il passe de 85 à
13. La mémorisation des états évite ici des développements répétés. L'ampleur
de l'amélioration varie néanmoins selon la stratégie et l'instance : sur le
labyrinthe simple, Greedy et \AStar{} développent chacun 4 noeuds dans les
deux modes.

\subsection{Coût de solution et effort}

Sur le trajet Gatineau--Baie-Comeau, \UCS{} et \AStar{} retournent un trajet
de 881,9\,km en 9 étapes, contre 912,7\,km en 6 étapes pour \BFS{} et Greedy.
La solution comportant le moins d'étapes n'est donc pas nécessairement celle
de coût minimal. Sur le 8-Puzzle moyen, Greedy illustre le même compromis :
385 développements pour une solution de 48 coups, au lieu des 1\,440
développements et 28 coups d'\AStar{}.

Le labyrinthe pondéré de ce benchmark constitue une limite importante : toutes
les stratégies y trouvent un chemin de coût 9 et de profondeur 6. Il ne
fournit donc pas un exemple où le coût de chemin distingue \BFS{} de \UCS{}.
Une instance pondérée supplémentaire, conçue pour offrir des chemins de
longueurs et de coûts différents, serait nécessaire pour mesurer directement
ce phénomène.

\subsection{Résolution des cas impossibles}

Pour le jeu arithmétique avec les nombres $1,1,1,1$ et la cible 99, chaque
stratégie termine l'exploration après 97 noeuds générés et 20 noeuds expansés,
sans solution. Il s'agit d'un échec après épuisement de l'espace d'états, et
non d'un arrêt dû à une limite de temps ou de mémoire. Le CSV du benchmark ne
contient pas de colonne de statut détaillé permettant de distinguer ces états
pour d'autres scénarios; le résultat de ce cas a été contrôlé pendant
l'exécution.

\subsection{Limites de l'évaluation}

\begin{itemize}
    \item Le benchmark ciblé couvre les cas indiqués dans la section
    \cref{sec:experiences}; il ne faut pas le confondre avec un éventuel
    benchmark antérieur portant sur davantage d'instances.
    \item Les temps correspondent à une seule exécution par configuration et
    peuvent varier selon l'environnement. Ils ne constituent pas des
    estimations statistiques.
    \item \code{max\_open} mesure la taille maximale de la frontière, pas la
    consommation mémoire totale.
    \item L'expérience heuristique porte sur une seule instance de 8-Puzzle.
    Des répétitions sur plusieurs états initiaux seraient nécessaires pour
    généraliser la conclusion.
\end{itemize}
